\chapter{Multiplexing dell'I/O, server a eventi e pool di thread}

Due versioni della stessa chat (un processo per client, un thread per client) e un costo che nessuna delle due evita.
\begin{itemize}
    \item Entrambe pagano al momento della connessione e continuano a pagare: uno spazio di indirizzi o uno stack, per un
          client che tace quasi sempre.
    \item E abbiamo lasciato una questione aperta: una \texttt{write} verso un client lento congela tutta la chat, perché
          il lock della lista viene tenuto durante la scrittura.
    \item Qui il flusso di controllo smette di essere uno per client. \textbf{Uno solo osserva molte connessioni} e serve
          quelle pronte; un \textbf{insieme fisso di worker} è tenuto occupato da una coda.
    \item Poi l'altra direzione: un datagramma, molti destinatari, e nessuno che tiene una lista.
\end{itemize}

% ══════════════════════════════════════════════════════════════
\section{Aspettare su più di un descriptor}
% ══════════════════════════════════════════════════════════════

\begin{itemize}
    \item \texttt{read} aspetta su un descriptor. Mentre aspetta lì, tutto il resto che il programma potrebbe vedere gli è
          invisibile.
    \item È l'unica ragione per cui avevamo un flusso di controllo per connessione: il processo, o il thread, esisteva
          \textbf{per potersi bloccare}.
    \item \texttt{select} capovolge la cosa: ci si blocca una volta sola, \textbf{su un insieme di descriptor}, e si torna
          appena uno qualsiasi può essere servito senza bloccarsi.
    \item Che cosa compra: nessun thread per socket, nessun ciclo di polling che brucia il processore, e un solo posto in
          cui vive lo stato di ogni client.
    \item Che cosa non compra: il \textbf{parallelismo}. Un flusso di controllo fa ancora una cosa alla volta, su un core.
\end{itemize}

\dfn{Multiplexing dell'I/O}{
    Si parla di \textbf{multiplexing dell'I/O} quando un solo flusso di controllo attende \textbf{contemporaneamente} su
    più descriptor, e il kernel gli dice quali sono pronti. In POSIX la chiamata è \texttt{select}; \texttt{poll} ed
    \texttt{epoll} sono varianti della stessa idea.
}

\begin{lstlisting}[language={[POSIX]C}, style=wnumbers, title={select}]
#include <sys/select.h>

int select(int nfds, fd_set *readfds, fd_set *writefds, fd_set *exceptfds,
           struct timeval *timeout);
\end{lstlisting}

\begin{itemize}
    \item Si blocca finché almeno un descriptor degli insiemi può essere usato, oppure scade il timeout, oppure arriva un
          segnale.
    \item Restituisce $-1$ ed \texttt{errno} in caso di errore, 0 se il timeout è scaduto senza nulla di pronto, altrimenti
          \textbf{quanti descriptor sono pronti}, contati su tutti e tre gli insiemi.
    \item \texttt{nfds} non è un conteggio: è il \textbf{descriptor più alto più uno}, perché il kernel scandisce
          l'intervallo da 0 a \texttt{nfds-1}.
\end{itemize}

\subsection{Gli insiemi di descriptor}

Un \texttt{fd\_set} è una \textbf{maschera di bit}: un bit per ogni numero di descriptor, e quattro macro per toccarla.

\begin{table}[H]
\centering
\begin{tabular}{>{\ttfamily}l l}
\toprule
\normalfont\textbf{Macro} & \textbf{Che cosa fa} \\
\midrule
FD\_ZERO(fd\_set *set)         & svuota l'insieme: si comincia sempre da qui \\
FD\_SET(int fd, fd\_set *set)   & mette \texttt{fd} nell'insieme \\
FD\_CLR(int fd, fd\_set *set)   & toglie \texttt{fd} dall'insieme \\
FD\_ISSET(int fd, fd\_set *set) & \texttt{fd} è nell'insieme? Si usa dopo la chiamata \\
\bottomrule
\end{tabular}
\caption{Le quattro macro sugli insiemi di descriptor.}
\end{table}

\begin{itemize}
    \item \texttt{readfds} per i descriptor da cui si vuole leggere, \texttt{writefds} per quelli su cui si vuole scrivere;
          \texttt{NULL} significa non osservare nulla di quel tipo.
    \item La maschera ha dimensione fissa: \texttt{FD\_SETSIZE} vale 1024 su Linux, e un descriptor uguale o superiore non
          può proprio entrare in un insieme.
\end{itemize}

\subsection{Quando un descriptor è pronto}

\begin{table}[H]
\centering
\begin{tabular}{p{0.45\textwidth} p{0.45\textwidth}}
\toprule
\textbf{Pronto in lettura}: una \texttt{read} tornerebbe subito & \textbf{Pronto in scrittura}: una \texttt{write} non si
                                                                    bloccherebbe \\
\midrule
sono arrivati byte nel buffer di ricezione & c'è posto nel buffer di invio (\texttt{SO\_SNDLOWAT}, 2048 byte per
                                             default) e il socket è connesso \\
il peer ha chiuso il suo lato di scrittura, quindi \texttt{read} restituisce 0 & una \texttt{connect} non bloccante è
                                                                                  terminata, in un modo o nell'altro \\
è un socket di ascolto e c'è una connessione in attesa & il peer ha chiuso il suo lato di lettura: la \texttt{write}
                                                          torna con \texttt{EPIPE} \\
c'è un errore in sospeso sul socket & c'è un errore in sospeso sul socket \\
\bottomrule
\end{tabular}
\caption{Le condizioni che rendono pronto un socket.}
\end{table}

Pronto significa \textbf{non si bloccherà}, non \emph{riuscirà}: un errore è un'ottima ragione per essere pronti.
\texttt{exceptfds} è quasi sempre \texttt{NULL}: su Linux l'unica cosa che segnala sono i dati fuori banda.

\subsection{Il timeout}

\begin{lstlisting}[language={[POSIX]C}, style=wnumbers, title={Tre modi di aspettare}]
struct timeval {
    time_t      tv_sec;     /* secondi */
    suseconds_t tv_usec;    /* microsecondi, non millisecondi */
};

struct timeval tv   = { .tv_sec = 2, .tv_usec = 500000 };    /* 2,5 secondi */
struct timeval zero = { .tv_sec = 0, .tv_usec = 0 };

select(nfds, &rset, NULL, NULL, NULL);      /* per sempre */
select(nfds, &rset, NULL, NULL, &zero);     /* polling: torna subito */
select(nfds, &rset, NULL, NULL, &tv);       /* al massimo 2,5 s */
\end{lstlisting}

Su Linux \texttt{select} riscrive in quella struttura il tempo che \textbf{non} ha usato: dentro un ciclo va riempita di
nuovo prima di ogni chiamata, altrimenti l'attesa si riduce fino a niente.

\subsection{select consuma gli insiemi}

\begin{figure}[H]
\centering
\begin{tikzpicture}[every node/.style={font=\scriptsize\sffamily},
    on/.style={draw=netblue!70!black, fill=netblue!25, minimum width=1.0cm, minimum height=0.5cm, font=\scriptsize\ttfamily},
    rd/.style={draw=netgreen!70!black, fill=netgreen!35, minimum width=1.0cm, minimum height=0.5cm, font=\scriptsize\ttfamily},
    off/.style={draw=black!30, fill=white, text=black!35, minimum width=1.0cm, minimum height=0.5cm, font=\scriptsize\ttfamily},
    a/.style={-{Stealth[length=1.8mm]}, thick, black!60}]
  \node[anchor=east, font=\scriptsize\ttfamily] at (-0.8,0) {allset};
  \foreach \i in {3,...,7} \node[on] at ({(\i-3)*1.15},0) {fd \i};
  \node[anchor=west, text=black!60] at (5.5,0) {ciò che vogliamo osservare};
  \draw[a] (2.3,-0.3) -- node[right]{\texttt{rset = allset}} (2.3,-0.85);
  \node[anchor=east, font=\scriptsize\ttfamily] at (-0.8,-1.15) {rset};
  \foreach \i in {3,...,7} \node[on] at ({(\i-3)*1.15},-1.15) {fd \i};
  \node[anchor=west, text=black!60] at (5.5,-1.15) {ciò che entra};
  \draw[a] (2.3,-1.45) -- node[right]{\texttt{select(...)}} (2.3,-2.0);
  \node[anchor=east, font=\scriptsize\ttfamily] at (-0.8,-2.3) {rset};
  \node[off] at (0,-2.3)    {fd 3};
  \node[rd]  at (1.15,-2.3) {fd 4};
  \node[off] at (2.3,-2.3)  {fd 5};
  \node[off] at (3.45,-2.3) {fd 6};
  \node[rd]  at (4.6,-2.3)  {fd 7};
  \node[anchor=west, text=black!60] at (5.5,-2.3) {ciò che torna: 2};
\end{tikzpicture}
\caption{Gli insiemi sono argomenti \emph{e} risultati. I tre descriptor non pronti sono stati azzerati: un insieme
         costruito una volta, fuori dal ciclo, aspetta la seconda volta solo quelli pronti la prima, e dopo ancora niente.}
\end{figure}

\warning{L'insieme costruito una volta sola}{
    Si costruisce l'insieme una volta, prima del ciclo, e si chiama \texttt{select} nel ciclo. Che cosa aspetta la seconda
    chiamata? \textbf{Solo i descriptor che erano pronti la prima volta}. \texttt{select} sovrascrive l'insieme con il
    sottoinsieme pronto e azzera tutto il resto. La cura è tenere un insieme \emph{principale} e passarne una
    \textbf{copia} a ogni giro, oppure ricostruirlo da capo ogni volta.
}

% ══════════════════════════════════════════════════════════════
\section{Un client che ascolta due descriptor}
% ══════════════════════════════════════════════════════════════

\warning{Il client che non sente il server}{
    Il client legge una riga, la invia e legge la risposta. Il server però parla anche di sua iniziativa: una riga ogni
    dieci secondi. Quando la vede l'utente? \textbf{Quando scrive qualcosa}. Il client è fermo nella \texttt{read} sulla
    tastiera, e le righe del server si accumulano nel buffer del socket finché non torna.
    \begin{itemize}
        \item Non è un bug del server: il client può essere bloccato in uno solo dei due posti, e ha scelto quello
              sbagliato.
        \item Un secondo thread lo risolverebbe: di nuovo uno per descriptor, con uno schermo condiviso su cui scrivere.
    \end{itemize}
}

Per due descriptor basta \textbf{una chiamata}.

\begin{lstlisting}[language={[POSIX]C}, style=wnumbers, title={Osservare la tastiera e il socket}]
static void chat_client(int fd) {
    char buf[256];
    fd_set rset;
    int nfds = (fd > STDIN_FILENO ? fd : STDIN_FILENO) + 1;

    for (;;) {
        FD_ZERO(&rset);                             /* ricostruito a ogni giro */
        FD_SET(STDIN_FILENO, &rset);
        FD_SET(fd, &rset);

        if (select(nfds, &rset, NULL, NULL, NULL) < 0) {
            if (errno == EINTR) continue;           /* un segnale, non un errore */
            perror("select"); return;
        }
        if (FD_ISSET(STDIN_FILENO, &rset)) {        /* l'utente ha scritto */
            ssize_t n = read(STDIN_FILENO, buf, sizeof buf);
            if (n <= 0) return;                     /* fine dell'input: abbiamo finito */
            if (write(fd, buf, n) != n) { perror("write"); return; }
        }
        if (FD_ISSET(fd, &rset)) {                  /* ha parlato il server */
            ssize_t n = read(fd, buf, sizeof buf);
            if (n == 0) { fprintf(stderr, "server closed\n"); return; }
            if (n < 0) { perror("read"); return; }
            write(STDOUT_FILENO, buf, n);
        }
    }
}
\end{lstlisting}

\begin{itemize}
    \item Due descriptor, un'attesa: il client ora è bloccato \textbf{in entrambi i posti insieme}, ed è tutto il punto.
    \item Due \texttt{if}, non un \texttt{if/else}: \texttt{select} può segnalarli entrambi, e il secondo resterebbe in
          attesa fino alla prossima riga digitata.
\end{itemize}

% ══════════════════════════════════════════════════════════════
\section{Un solo flusso di controllo, molte connessioni}
% ══════════════════════════════════════════════════════════════

Lo stesso server, con il ciclo rovesciato: non aspetta più un client, aspetta \textbf{quello che parla per primo}.

\begin{figure}[H]
\centering
\begin{tikzpicture}[every node/.style={font=\scriptsize\sffamily},
    st/.style={draw=netblue!70!black, fill=netblue!8, rounded corners=2pt, align=center, minimum width=3.0cm,
               minimum height=1.3cm},
    a/.style={-{Stealth[length=2mm]}, thick}]
  \node[font=\small\bfseries\sffamily] at (4,1.25) {un processo, un thread, nessun lock};
  \node[st] (w) at (0,0) {\textbf{l'insieme osservato}\\\texttt{fd 3} ascolto\\\texttt{fd 4 5 6 7} client};
  \node[st, draw=netred!70!black, fill=netred!6] (s) at (4,0) {\texttt{\textbf{select(...)}}\\l'unica chiamata\\che si blocca};
  \node[st, draw=netgreen!70!black, fill=netgreen!8] (r) at (8,0) {\textbf{i pronti}\\\texttt{fd 4} e \texttt{fd 7}\\e nient'altro};
  \draw[a] (w) -- (s);
  \draw[a] (s) -- (r);
  \draw[a] (r.south) -- ++(0,-0.6) -| node[pos=0.25, below]{\texttt{accept}, \texttt{read}, \texttt{write}, e di nuovo} (w.south);
\end{tikzpicture}
\caption{Ogni client è un posto in un array e un bit in un insieme. Non c'è nulla da proteggere con un lock, perché nient'altro
         è in esecuzione; e nulla sulla via del ritorno può aspettare, perché tutti aspetterebbero con lui.}
\end{figure}

\dfn{Server guidato dagli eventi}{
    Un server \textbf{guidato dagli eventi} (\emph{event-driven}) ha un solo ciclo, il \key{ciclo degli eventi}: si blocca
    in \texttt{select}, poi esegue per ogni descriptor pronto un pezzo di lavoro breve e che non si blocca mai, e torna in
    \texttt{select}. Lo stato di ogni client sta in strutture dati, non sullo stack di un thread.
}

\begin{lstlisting}[language={[POSIX]C}, style=wnumbers, title={Lo stato del server}]
#define MAXC 64
static int clients[MAXC];                    /* socket connessi, -1 = posto libero */

int main(int argc, char *argv[]) {
    int listen_fd = listening_socket(argc > 1 ? argv[1] : "5000");
    fd_set allset, rset;                     /* l'insieme principale, e la copia */
    int maxfd = listen_fd;

    for (int i = 0; i < MAXC; i++) clients[i] = -1;
    FD_ZERO(&allset);
    FD_SET(listen_fd, &allset);              /* la porta e' un descriptor come gli altri */

    for (;;) {
        rset = allset;                       /* assegnamento di struct: una copia */
        int ready = select(maxfd + 1, &rset, NULL, NULL, NULL);
        if (ready < 0) { if (errno == EINTR) continue; perror("select"); break; }
\end{lstlisting}

La stessa lista di client della versione con i thread, \textbf{senza il mutex}: non c'è nient'altro in esecuzione.

\begin{lstlisting}[language={[POSIX]C}, style=wnumbers, title={Arriva una connessione sul socket di ascolto}]
        if (FD_ISSET(listen_fd, &rset)) {
            int fd = accept(listen_fd, NULL, NULL);
            if (fd < 0) { perror("accept"); continue; }

            int slot = -1;
            for (int i = 0; i < MAXC; i++)
                if (clients[i] == -1) { clients[i] = fd; slot = i; break; }

            if (slot < 0) { write(fd, "chat is full\n", 13); close(fd); }
            else {
                FD_SET(fd, &allset);         /* osservato dal prossimo giro */
                if (fd > maxfd) maxfd = fd;
            }
            if (--ready == 0) continue;      /* nient'altro era pronto */
        }
\end{lstlisting}

Un socket di ascolto leggibile significa una connessione in attesa, quindi questa \texttt{accept} ritorna subito.

\begin{lstlisting}[language={[POSIX]C}, style=wnumbers, title={Servire le connessioni pronte}]
        for (int i = 0; i < MAXC && ready > 0; i++) {
            int fd = clients[i];
            if (fd == -1 || !FD_ISSET(fd, &rset)) continue;
            --ready;

            char buf[256];
            ssize_t n = read(fd, buf, sizeof buf);
            if (n > 0) {
                broadcast(i, buf, n);        /* tutti tranne i */
            } else {                         /* 0: se n'e' andato. <0: si e' rotto */
                FD_CLR(fd, &allset);
                close(fd);
                clients[i] = -1;
            }
        }
    }
}
\end{lstlisting}

\begin{itemize}
    \item \texttt{broadcast} è il \texttt{for} della versione con i thread, meno il mutex.
    \item \texttt{maxfd} non si abbassa quando un client se ne va: costa una scansione un po' più lunga, e l'alternativa è
          una ricerca a ogni uscita.
    \item \texttt{ready} serve a fermarsi presto: quando sono stati serviti tanti descriptor quanti ne ha contati
          \texttt{select}, il resto dell'array non va guardato.
\end{itemize}

\subsection{In questo ciclo nulla può bloccarsi}

\warning{Mezza riga, un byte al secondo}{
    Un client invia mezza riga, un byte alla volta, un byte al secondo. La \texttt{read} ha restituito sei byte e la riga
    non è finita. E adesso? Si prendono i sei byte, si mettono \textbf{nel buffer di quel client} e si torna in
    \texttt{select}. Aspettare il resto vorrebbe dire aspettare un client con il tempo di tutti gli altri.
}

\begin{itemize}
    \item Ogni client quindi ha bisogno del proprio \textbf{buffer di ingresso}: il ciclo tiene lo stato che un thread teneva
          sul suo stack.
    \item E \emph{pronto} non è una promessa: fra la \texttt{select} e l'\texttt{accept} il client può mandare un reset, e
          l'\texttt{accept} si blocca, con tutto il server dentro.
    \item Da qui i descriptor \textbf{non bloccanti}: \texttt{EAGAIN} invece di un'attesa, e il ciclo va avanti.
\end{itemize}

\begin{lstlisting}[language={[POSIX]C}, style=wnumbers, title={Un descriptor non bloccante}]
int flags = fcntl(fd, F_GETFL, 0);
fcntl(fd, F_SETFL, flags | O_NONBLOCK);      /* da ora: mai un'attesa */

ssize_t n = read(fd, buf, sizeof buf);
if (n < 0 && (errno == EAGAIN || errno == EWOULDBLOCK))
    ;                                        /* niente per ora: si torna in select */

n = recv(fd, buf, sizeof buf, MSG_DONTWAIT); /* lo stesso, per una sola chiamata */
\end{lstlisting}

\subsection{Il client lento, finalmente}

\begin{lstlisting}[language={[POSIX]C}, style=wnumbers, title={Una coda in uscita per client}]
static struct { char buf[4096]; size_t len; } out[MAXC];   /* byte in attesa, per client */
static fd_set wall;                          /* chi ha ancora qualcosa da inviare */

static void queue(int i, const char *msg, size_t n) {      /* non si blocca mai */
    if (out[i].len + n > sizeof out[i].buf) { drop(i); return; }  /* troppo lento: via */
    memcpy(out[i].buf + out[i].len, msg, n);
    out[i].len += n;
    FD_SET(clients[i], &wall);               /* avvisami quando puo' ricevere byte */
}

/* e nel ciclo, per ogni i che select ha segnalato scrivibile */
ssize_t w = write(clients[i], out[i].buf, out[i].len);
if (w > 0) { out[i].len -= w; memmove(out[i].buf, out[i].buf + w, out[i].len); }
if (out[i].len == 0) FD_CLR(clients[i], &wall);            /* niente da inviare: basta chiedere */
\end{lstlisting}

Ora \texttt{select} riceve anche un insieme di scrittura. \texttt{broadcast} non scrive più: \textbf{accoda}. Il client
lento riempie il proprio buffer e viene scartato quando trabocca, e nessun altro lo aspetta.

\warning{Perché un descriptor entra in wall solo se ha qualcosa}{
    Un socket connesso è scrivibile quasi sempre, perché il suo buffer di invio è quasi sempre vuoto. Se tutti i client
    stessero sempre nell'insieme di scrittura, \texttt{select} tornerebbe subito a ogni giro: il server girerebbe al
    100\% di CPU senza fare nulla. Un descriptor entra nell'insieme di scrittura quando la sua coda si riempie e ne esce
    quando si svuota.
}

\subsection{Quanto costa select}

\begin{itemize}
    \item Gli insiemi sono copiati nel kernel a ogni chiamata, e il kernel percorre tutto l'intervallo da 0 a
          \texttt{nfds-1}: il costo segue il \textbf{descriptor più alto}, non il numero di quelli pronti.
    \item Al ritorno si percorre il proprio array per trovare chi è pronto: un secondo passaggio lineare su client che per lo
          più non hanno detto niente.
    \item E \texttt{FD\_SETSIZE} vale 1024: un server con diecimila connessioni non può nemmeno nominare i propri descriptor
          in un insieme.
    \item \texttt{poll} prende un array invece di una maschera: nessun limite fisso, ma ancora lineare. \texttt{epoll} su
          Linux e \texttt{kqueue} sui BSD \textbf{registrano i descriptor una volta sola} e restituiscono solo quelli
          pronti.
\end{itemize}

\begin{lstlisting}[language={[POSIX]C}, style=wnumbers, title={Lo stesso ciclo con epoll}]
int ep = epoll_create1(0);
struct epoll_event ev = { .events = EPOLLIN, .data.fd = listen_fd };
epoll_ctl(ep, EPOLL_CTL_ADD, listen_fd, &ev);          /* registrato una volta */

struct epoll_event ready[64];
for (;;) {
    int n = epoll_wait(ep, ready, 64, -1);             /* -1: per sempre */
    for (int k = 0; k < n; k++) {                      /* solo i pronti, nessuna scansione */
        int fd = ready[k].data.fd;
        if (fd == listen_fd) { /* accept, poi EPOLL_CTL_ADD del nuovo fd */ }
        else                 { /* read, oppure EPOLL_CTL_DEL e close */ }
    }
}
\end{lstlisting}

La forma del ciclo non cambia: cambia la chiamata nel mezzo. Tutto quello che vale per \texttt{select} vale ancora.

% ══════════════════════════════════════════════════════════════
\section{Un insieme fisso di worker}
% ══════════════════════════════════════════════════════════════

Il ciclo degli eventi ha eliminato il flusso di controllo per client e ne ha lasciato uno per tutti. Un \textbf{pool} ne
tiene diversi, ma in numero fisso, creati prima che arrivi il primo client.

\dfn{Pool di thread}{
    Un \textbf{pool di thread} è un insieme di $n$ thread creati all'avvio, che prendono il lavoro da una \textbf{coda
    limitata} e, finito un lavoro, tornano a prenderne un altro. Chi produce il lavoro lo accoda e basta.
}

\begin{figure}[H]
\centering
\begin{tikzpicture}[every node/.style={font=\scriptsize\sffamily},
    cell/.style={draw=black!60, minimum width=0.55cm, minimum height=0.55cm, font=\scriptsize\ttfamily, fill=#1},
    wk/.style={draw=netgreen!70!black, fill=netgreen!8, rounded corners=2pt, minimum width=3.5cm, minimum height=0.55cm,
               anchor=west},
    a/.style={-{Stealth[length=2mm]}, thick}]
  \node[draw=netblue!70!black, fill=netblue!8, rounded corners=2pt, align=center, minimum height=1.2cm] (acc) at (0,0)
    {\texttt{\textbf{accept}}\\un ciclo\\nient'altro};
  \foreach \i/\v/\c in {0//white, 1//white, 2/4/netblue!20, 3/7/netblue!20, 4/9/netblue!20}
    \node[cell=\c] (q\i) at (2.6+\i*0.55,0) {\v};
  \node[above=0.25cm of q2, text=black!70] {una coda limitata di descriptor};
  \draw[a] (acc.east) -- node[above]{\texttt{put\_fd}} (q0.west);
  \node[wk] (w1) at (6.7,0.8)  {worker 1: \texttt{serve(fd 4)}};
  \node[wk] (w2) at (6.7,0)    {worker 2: \texttt{serve(fd 7)}};
  \node[wk, text=black!50] (w3) at (6.7,-0.8) {worker 3: \texttt{get\_fd}, dorme};
  \node[above=0.1cm of w1, text=black!70] {creati una volta, prima di ogni client};
  \draw[thick] (q4.east) -- node[above]{\texttt{get\_fd}} (6.0,0);
  \draw[a] (6.0,0) -- (w2.west);
  \draw[a] (6.0,0) -- (w1.west);
  \draw[a] (6.0,0) -- (w3.west);
  \node[text=netred, align=center] at (0,-1.15) {coda piena:\\questo ciclo aspetta};
  \draw[-{Stealth[length=1.8mm]}, dashed, black!55] (w3.south) -- ++(0,-0.45) -| node[pos=0.25, below]{\texttt{close},
    poi di nuovo a prendere il prossimo} (q2.south);
\end{tikzpicture}
\caption{Una connessione costa un inserimento in coda e un risveglio: quando arriva un client non si crea niente, e niente
         si distrugge.}
\end{figure}

\begin{lstlisting}[language={[POSIX]C}, style=wnumbers, title={La coda è il buffer limitato, con dentro descriptor}]
#define QSIZE 128

typedef struct {
    int fd[QSIZE];                           /* descriptor invece di interi */
    int in, out, count;
    pthread_mutex_t mtx;
    pthread_cond_t not_full, not_empty;      /* le stesse due attese di prima */
} queue_t;

void put_fd(queue_t *q, int fd);             /* si blocca finche' la coda e' piena */
int  get_fd(queue_t *q);                     /* si blocca finche' e' vuota */
\end{lstlisting}

È il buffer produttore-consumatore, intatto: un mutex sullo stato, una variabile di condizione per ogni cosa che qualcuno
aspetta. Quello che lo attraversa ora è una connessione.

\begin{lstlisting}[language={[POSIX]C}, style=wnumbers, title={Il ciclo di accept è il produttore}]
int main(int argc, char *argv[]) {
    int listen_fd = listening_socket(argc > 1 ? argv[1] : "5000");
    queue_t q; queue_init(&q);
    pthread_t tid;

    for (int i = 0; i < NWORKERS; i++)       /* una volta, prima di ogni client */
        if (pthread_create(&tid, NULL, worker, &q) == 0) pthread_detach(tid);

    for (;;) {
        int fd = accept(listen_fd, NULL, NULL);
        if (fd < 0) { if (errno == EINTR) continue; perror("accept"); continue; }
        put_fd(&q, fd);                      /* l'unica cosa che fa il ciclo */
    }
}
\end{lstlisting}

Niente \texttt{fork}, niente \texttt{pthread\_create} nel ciclo, niente \texttt{malloc} per un argomento: una connessione
costa un inserimento in coda.

\begin{lstlisting}[language={[POSIX]C}, style=wnumbers, title={Il worker è il consumatore}]
static void *worker(void *arg) {
    queue_t *q = arg;

    for (;;) {
        int fd = get_fd(q);                  /* dorme qui finche' non c'e' niente */
        if (fd == STOP) break;               /* la pillola avvelenata: una per worker */

        serve(fd);                           /* la funzione del server sequenziale */
        close(fd);
    }
    return NULL;
}
\end{lstlisting}

Il thread c'era prima del client e c'è dopo: per connessione non si crea né si distrugge niente. E \texttt{serve()} è
passata ormai per tre cicli diversi senza essere toccata.

\subsection{Quanti worker, e che cosa succede quando la coda si riempie}

\begin{itemize}
    \item Il numero non è «quanti sono i client»: è \textbf{quanti possono lavorare insieme}. Circa quanti sono i core per un
          lavoro che calcola, di più quando i worker passano il tempo bloccati in \texttt{read}.
    \item La coda è limitata \textbf{di proposito}. Quando si riempie, \texttt{put\_fd} blocca il ciclo di \texttt{accept};
          le connessioni si accumulano allora nella coda di listen, e quando si riempie anche quella il kernel le rifiuta.
    \item Una coda illimitata sposta solo il guasto più in profondità: accetta tutti e non risponde a nessuno, e il client
          ha rinunciato molto prima che arrivi il suo turno.
\end{itemize}

\dfn{Contropressione}{
    La catena coda piena $\to$ \texttt{accept} ferma $\to$ coda di listen piena $\to$ connessioni rifiutate è la
    \textbf{contropressione} (\emph{back pressure}): un server che non riesce a stare al passo lo dice, alla porta, invece
    di allocare memoria finché non muore.
}

Il numero da misurare quindi non è il throughput, ma \textbf{quanto aspetta una richiesta} prima che un worker la prenda.

\warning{Tutti e sedici i worker parlano con qualcuno}{
    Sedici worker, sedici client di chat connessi e silenziosi. Si connette il diciassettesimo: quando viene servito?
    \textbf{Mai}. Ogni worker è dentro \texttt{serve()}, bloccato in \texttt{read} sul proprio client, e \texttt{serve()}
    ritorna solo quando quel client se ne va. La coda non è piena: semplicemente non viene mai letta.
    \begin{itemize}
        \item Un pool il cui lavoro è una connessione serve tanti client quanti worker ha. È la forma giusta quando il
              lavoro \textbf{finisce}: una richiesta HTTP, una query, un file da inviare.
        \item Per connessioni che durano, o si tiene un flusso per client limitandone il numero (il semaforo contatore già
              visto), oppure il lavoro diventa \textbf{il messaggio}.
        \item Il che richiede qualcuno che si accorga che un messaggio è arrivato: il ciclo di prima. Uno che osserva, un
              pool che lavora.
    \end{itemize}
}

\subsection{Tre forme di server}

\begin{table}[H]
\centering
\begin{tabular}{l l l l}
\toprule
 & \textbf{Uno per client} & \textbf{Un ciclo} & \textbf{Pool fisso} \\
\midrule
costo di un nuovo client  & uno stack          & un posto e un bit      & un inserimento in coda \\
quanti insieme            & migliaia           & decine di migliaia     & quanti sono i worker \\
una chiamata bloccante    & ferma quel client  & ferma tutti            & occupa un worker \\
stato condiviso           & vuole un mutex     & non vuole niente       & vuole un mutex \\
usa più core              & sì                 & no                     & sì \\
\bottomrule
\end{tabular}
\caption{Le tre forme di server a confronto.}
\end{table}

Non restano alternative a lungo: un server reale esegue \textbf{un ciclo per core}, ciascuno con un pool alle spalle per
il lavoro che non deve fermare il ciclo.

% ══════════════════════════════════════════════════════════════
\section{Un mittente, molti destinatari}
% ══════════════════════════════════════════════════════════════

\begin{figure}[H]
\centering
\begin{tikzpicture}[every node/.style={font=\scriptsize\sffamily},
    a/.style={-{Stealth[length=1.8mm]}, thick}]
  % broadcast
  \node[font=\small\bfseries\sffamily] at (2,2.3) {broadcast \texttt{192.168.1.255}};
  \node (s1) at (2,1.45) {\icon[0.75cm]{pc}};
  \draw[wired] (0,0.5) -- (4,0.5);
  \draw[wired] (s1.south) -- (2,0.5);
  \foreach \x in {0.2,1.4,2.6,3.8} {
    \node (h) at (\x,-0.45) {\icon[0.65cm]{pc}};
    \draw[a, netgreen!70!black] (\x,0.5) -- (h.north);
  }
  \node[align=center, text=black!70] at (2,-1.6) {la riceve ogni host della sottorete:\\nessuno l'ha chiesta, nessuno può
    rifiutarla};
  % multicast
  \node[font=\small\bfseries\sffamily] at (8,2.3) {multicast \texttt{239.1.2.3}};
  \node (s2) at (8,1.45) {\icon[0.75cm]{pc}};
  \draw[wired] (6,0.5) -- (10,0.5);
  \draw[wired] (s2.south) -- (8,0.5);
  \node[opacity=0.35] at (6.2,-0.45) {\icon[0.65cm]{pc}};
  \node (j1) at (7.4,-0.45) {\icon[0.65cm]{pc}};
  \node[opacity=0.35] at (8.6,-0.45) {\icon[0.65cm]{pc}};
  \node (j2) at (9.8,-0.45) {\icon[0.65cm]{pc}};
  \draw[a, netgreen!70!black] (7.4,0.5) -- (j1.north);
  \draw[a, netgreen!70!black] (9.8,0.5) -- (j2.north);
  \node[below=-2pt of j1, text=netgreen!60!black] {iscritto};
  \node[below=-2pt of j2, text=netgreen!60!black] {iscritto};
  \node[align=center, text=black!70] at (8,-1.6) {solo gli host iscritti al gruppo; una copia\\sul filo, qualunque sia il
    numero di membri};
\end{tikzpicture}
\caption{Entrambi sono datagrammi: non esiste una connessione con molti peer, quindi niente TCP. E solo IPv4: IPv6 ha
         eliminato il broadcast e ne affida i compiti al multicast.}
\end{figure}

\begin{lstlisting}[language={[POSIX]C}, style=wnumbers, title={Broadcast: i due estremi}]
/* il mittente */
int fd = socket(AF_INET, SOCK_DGRAM, 0);
int yes = 1;
setsockopt(fd, SOL_SOCKET, SO_BROADCAST, &yes, sizeof yes);   /* spento per default */
to.sin_family = AF_INET;
to.sin_port   = htons(5000);
inet_pton(AF_INET, "192.168.1.255", &to.sin_addr);  /* il broadcast di questa sottorete */
sendto(fd, msg, len, 0, (struct sockaddr *) &to, sizeof to);

/* il ricevente: un normale server UDP */
setsockopt(fd, SOL_SOCKET, SO_REUSEADDR, &yes, sizeof yes);
me.sin_family      = AF_INET;
me.sin_addr.s_addr = htonl(INADDR_ANY);             /* non l'indirizzo di broadcast */
me.sin_port        = htons(5000);                   /* la stessa porta del mittente */
bind(fd, (struct sockaddr *) &me, sizeof me);
recvfrom(fd, buf, sizeof buf, 0, (struct sockaddr *) &from, &flen);
\end{lstlisting}

\texttt{SO\_BROADCAST} esiste perché nessuno inondi una LAN sbagliando a scrivere un indirizzo: bisogna dire che lo si
voleva davvero. Senza, la \texttt{sendto} verso un indirizzo di broadcast fallisce con \texttt{EACCES}.

\warning{Il broadcast che non arriva}{
    Il programma gira, \texttt{sendto} restituisce il numero di byte inviati, e nessuno riceve niente. Dove è finito?
    \textbf{Fuori dal socket, e non oltre}. \texttt{sendto} riferisce ciò che ha lasciato la nostra macchina: con UDP è
    tutto ciò che può promettere.
    \begin{itemize}
        \item \texttt{255.255.255.255} è il \textbf{broadcast limitato}: non lascia mai l'interfaccia da cui è inviato. I
              router non inoltrano affatto il broadcast: si ferma alla sottorete, per progetto.
        \item Il broadcast di sottorete deve essere \textbf{quello di quella sottorete}: va letto, non indovinato.
        \item Una macchina virtuale in modalità NAT sta su una rete privata tutta sua, e quel NAT non inoltra il broadcast.
              In modalità bridge prende un indirizzo sulla LAN vera e l'aula lo riceve, se il driver del Wi-Fi permette il
              bridging.
    \end{itemize}
}

\begin{lstlisting}[language=bash, title={L'indirizzo di broadcast, letto e non indovinato}]
$ ip addr show ens33 | grep inet
    inet 192.168.1.23/24 brd 192.168.1.255 scope global ens33     <- brd: quello giusto
\end{lstlisting}

\subsection{Gli indirizzi multicast}

\dfn{Gruppo multicast}{
    Un \textbf{gruppo multicast} è un indirizzo IPv4 fra \texttt{224.0.0.0} e \texttt{239.255.255.255}. Nomina
    \textbf{una conversazione, non un host}: nessuno lo possiede, e qualsiasi host può iscriversi.
}

\begin{table}[H]
\centering
\begin{tabular}{>{\ttfamily}l l}
\toprule
\normalfont\textbf{Intervallo} & \textbf{Uso} \\
\midrule
224.0.0.0/24              & link-local: mai inoltrato; qui vivono i protocolli di routing \\
224.0.1.0 -- 238.255.255.255 & ambito globale: richiede router configurati per questo \\
239.0.0.0/8               & ambito amministrativo: locale a un'organizzazione \\
\bottomrule
\end{tabular}
\caption{Gli ambiti degli indirizzi multicast IPv4.}
\end{table}

Per qualsiasi cosa nostra si usa \texttt{239.x.y.z}. L'altro limite a quanto lontano viaggia un datagramma è il suo
\textbf{TTL}, che per il multicast il mittente sceglie esplicitamente.

\subsection{Entrare in un gruppo}

\begin{lstlisting}[language={[POSIX]C}, style=wnumbers, title={Iscriversi a 239.1.2.3 sulla porta 5000}]
int fd = socket(AF_INET, SOCK_DGRAM, 0);
int yes = 1;
setsockopt(fd, SOL_SOCKET, SO_REUSEADDR, &yes, sizeof yes);
setsockopt(fd, SOL_SOCKET, SO_REUSEPORT, &yes, sizeof yes);

me.sin_family      = AF_INET;
me.sin_addr.s_addr = htonl(INADDR_ANY);
me.sin_port        = htons(5000);
bind(fd, (struct sockaddr *) &me, sizeof me);         /* la porta, non il gruppo */

struct ip_mreq mreq;
inet_pton(AF_INET, "239.1.2.3", &mreq.imr_multiaddr); /* quale gruppo */
mreq.imr_interface.s_addr = htonl(INADDR_ANY);        /* su quale interfaccia */
setsockopt(fd, IPPROTO_IP, IP_ADD_MEMBERSHIP, &mreq, sizeof mreq);
\end{lstlisting}

\begin{itemize}
    \item \texttt{bind} sceglie la porta, l'iscrizione sceglie il gruppo: senza di essa il kernel butta via i datagrammi
          anche quando arrivano sul filo.
    \item \texttt{SO\_REUSEADDR} riusa una porta lasciata da un socket precedente; \texttt{SO\_REUSEPORT} permette a più
          socket vivi di tenere la stessa porta nello stesso momento, ed è ciò che serve per eseguire due peer sulla stessa
          macchina.
\end{itemize}

\subsection{Inviare a un gruppo}

\begin{lstlisting}[language={[POSIX]C}, style=wnumbers, title={Il mittente}]
int fd = socket(AF_INET, SOCK_DGRAM, 0);              /* un normale socket UDP */

unsigned char ttl = 1;                                /* 1: questa sottorete e non oltre */
setsockopt(fd, IPPROTO_IP, IP_MULTICAST_TTL, &ttl, sizeof ttl);

unsigned char loop = 1;                               /* anche ai membri su questo host */
setsockopt(fd, IPPROTO_IP, IP_MULTICAST_LOOP, &loop, sizeof loop);

to.sin_family = AF_INET;
to.sin_port   = htons(5000);
inet_pton(AF_INET, "239.1.2.3", &to.sin_addr);
sendto(fd, msg, len, 0, (struct sockaddr *) &to, sizeof to);
\end{lstlisting}

Il mittente non si iscrive a niente, non si connette a niente e non sa chi sta ascoltando: questo codice è identico che
il gruppo abbia zero membri o duecento. TTL 1 e loop attivo sono anche i valori predefiniti; qui sono scritti per
mostrarli.

\subsection{Chi rende selettivo il multicast}

\begin{figure}[H]
\centering
\begin{tikzpicture}[every node/.style={font=\scriptsize\sffamily},
    a/.style={-{Stealth[length=1.8mm]}, thick}]
  \node (h1) at (0,1.0)  {\icon[0.6cm]{pc}};
  \node (h2) at (0,0)    {\icon[0.6cm]{pc}};
  \node[opacity=0.35] (h3) at (0,-1.0) {\icon[0.6cm]{pc}};
  \node[left=0pt of h1, text=netgreen!60!black] {iscritto};
  \node[left=0pt of h2, text=netgreen!60!black] {iscritto};
  \node[left=0pt of h3, text=black!50] {--};
  \node (sw) at (2.6,0) {\icon[0.9cm]{switch}};
  \node[below=0pt of sw, align=center] {switch\\IGMP snooping};
  \node (ra) at (5.4,0) {\icon[0.9cm]{router}};
  \node[below=0pt of ra, align=center] {router A\\questa sottorete};
  \node (rb) at (8.2,0) {\icon[0.9cm]{router}};
  \node[below=0pt of rb] {router B};
  \node (src) at (10.6,0) {\icon[0.75cm]{server}};
  \node[below=0pt of src, align=center] {sorgente\\del gruppo};
  \draw[wired] (h1.east) -- (sw.west);
  \draw[wired] (h2.east) -- (sw.west);
  \draw[wired, black!30] (h3.east) -- (sw.west);
  \draw[wired] (sw.east) -- (ra.west);
  \draw[wired] (ra.east) -- node[above]{PIM} (rb.west);
  \draw[wired] (rb.east) -- node[above]{PIM} (src.west);
  \draw[a, netred, dashed] (h1.north east) to[out=30, in=150] node[above]{IGMP report} (ra.north west);
  \node[text=black!60, align=center] at (2.6,-1.6) {solo le porte\\che l'hanno chiesto};
  \node[text=black!60, align=center] at (6.8,-1.6) {l'albero raggiunge solo le\\sottoreti che hanno un membro};
\end{tikzpicture}
\caption{I tre meccanismi che portano un datagramma multicast solo dove serve.}
\end{figure}

\begin{itemize}
    \item \key{IGMP}: l'iscrizione non è solo locale. Il kernel annuncia il gruppo al router della sottorete, e lo ripete
          quando gli viene chiesto.
    \item \key{PIM}: i router costruiscono un albero fra sottoreti, così che una sottorete senza membri non riceva niente.
    \item \key{IGMP snooping}: lo switch inoltra il gruppo solo sulle porte che l'hanno chiesto; senza, lo inonda su tutta
          la LAN e il filtro lo fa ciascun kernel.
    \item Nessuno dei tre è un controllo di accesso: qualsiasi host lungo il percorso può iscriversi e leggere tutto. Se
          deve restare privato, \textbf{si cifra il contenuto}.
\end{itemize}

% ══════════════════════════════════════════════════════════════
\section{Esercizi}
% ══════════════════════════════════════════════════════════════

\begin{esercizio}{fork, pipe e select}
Un genitore e un figlio osservano entrambi la tastiera con \texttt{select}; il genitore inoltra al figlio, attraverso una
pipe, ciò che legge. Il figlio osserva sia la tastiera sia la pipe e stampa da dove arriva ogni riga. Scrivere qualche
stringa e spiegare che cosa si vede: chi riceve la riga scritta, e si può dire in anticipo?

\svolgimento
\begin{lstlisting}[language={[POSIX]C}, style=wnumbers, title={fork\_pipe\_select.c}]
int p[2]; char buf[256]; ssize_t n;
pipe(p);
if (fork() == 0) {                           /* figlio: tastiera e pipe */
    close(p[1]);
    for (;;) {
        fd_set r; FD_ZERO(&r); FD_SET(STDIN_FILENO, &r); FD_SET(p[0], &r);
        if (select(p[0] + 1, &r, NULL, NULL, NULL) < 0) continue;
        if (FD_ISSET(STDIN_FILENO, &r) && (n = read(STDIN_FILENO, buf, sizeof buf)) > 0)
            printf("figlio, dalla tastiera: %.*s", (int) n, buf);
        if (FD_ISSET(p[0], &r) && (n = read(p[0], buf, sizeof buf)) > 0)
            printf("figlio, dal genitore:   %.*s", (int) n, buf);
    }
}
close(p[0]);                                 /* genitore: tastiera, poi pipe */
for (;;) {
    fd_set r; FD_ZERO(&r); FD_SET(STDIN_FILENO, &r);
    if (select(STDIN_FILENO + 1, &r, NULL, NULL, NULL) < 0) continue;
    if ((n = read(STDIN_FILENO, buf, sizeof buf)) <= 0) break;
    write(p[1], buf, n);
}
\end{lstlisting}

I due processi condividono lo \textbf{stesso} standard input, ereditato con la \texttt{fork}. Quando si preme Invio,
\texttt{select} sveglia \textbf{entrambi}: per tutti e due la tastiera è pronta. Ma la riga è una sola, e la prende chi
arriva prima alla \texttt{read}, cosa che decide lo scheduler. Non si può dire in anticipo.

L'altro processo, svegliato per niente, chiama \texttt{read} e \textbf{si blocca} fino alla riga successiva. Se è il
figlio, la riga che il genitore gli ha appena inoltrato resta nella pipe finché non si scrive ancora qualcosa: le righe
\emph{dal genitore} compaiono in ritardo e fuori ordine. È il «pronto non è una promessa» del capitolo, con due processi:
fra la \texttt{select} e la \texttt{read} qualcun altro ha preso i dati. Mettere lo standard input in
\texttt{O\_NONBLOCK} lo trasforma in un \texttt{EAGAIN}; ma il flag sta nella descrizione del file aperto, quindi vale per
entrambi i processi, e anche per la shell dopo che il programma è finito.
\end{esercizio}

\begin{esercizio}{Un peer multicast}
Scrivere un peer: un solo socket UDP, iscritto a \texttt{239.1.2.3:5000}, con \texttt{select} sullo standard input e sul
socket. Ciò che si scrive va al gruppo, ciò che arriva va sullo schermo. Nessun server: ogni peer è entrambe le cose.
Provarne tre sulla stessa macchina (\texttt{./peer 239.1.2.3 5000}), poi su due macchine. Perché si vedono i propri
messaggi, e quali sono i due modi per non vederli più?

\svolgimento
\begin{lstlisting}[language={[POSIX]C}, style=wnumbers, title={peer.c: il cuore}]
/* socket, SO_REUSEADDR, SO_REUSEPORT, bind sulla porta e IP_ADD_MEMBERSHIP come nel capitolo */
struct sockaddr_in to = me;                  /* stessa porta ... */
inet_pton(AF_INET, argv[1], &to.sin_addr);   /* ... ma l'indirizzo del gruppo */

for (;;) {
    fd_set r; FD_ZERO(&r); FD_SET(STDIN_FILENO, &r); FD_SET(fd, &r);
    if (select(fd + 1, &r, NULL, NULL, NULL) < 0) { if (errno == EINTR) continue; break; }

    if (FD_ISSET(STDIN_FILENO, &r)) {        /* da tastiera al gruppo */
        ssize_t n = read(STDIN_FILENO, buf, sizeof buf);
        if (n <= 0) break;
        sendto(fd, buf, n, 0, (struct sockaddr *) &to, sizeof to);
    }
    if (FD_ISSET(fd, &r)) {                  /* dal gruppo allo schermo */
        struct sockaddr_in from; socklen_t fl = sizeof from;
        char ip[INET_ADDRSTRLEN];
        ssize_t n = recvfrom(fd, buf, sizeof buf, 0, (struct sockaddr *) &from, &fl);
        inet_ntop(AF_INET, &from.sin_addr, ip, sizeof ip);
        printf("[%s:%d] %.*s", ip, ntohs(from.sin_port), (int) n, buf);
    }
}
\end{lstlisting}

Ci si vede perché \texttt{IP\_MULTICAST\_LOOP} vale 1 per default e il nostro socket è iscritto al gruppo: la copia
rimandata indietro sull'host arriva a tutti i socket iscritti, compreso il mittente. I due modi:
\begin{itemize}
    \item \textbf{Spegnere il loop} con \texttt{IP\_MULTICAST\_LOOP} a 0. Ma la copia locale è una sola per tutto l'host:
          spegnendola, non sentono più i nostri messaggi nemmeno gli altri peer sulla stessa macchina.
    \item \textbf{Riconoscere i propri messaggi} e scartarli. L'indirizzo del mittente non basta, perché sulla stessa
          macchina tutti i peer hanno lo stesso IP e la stessa porta 5000: serve un identificatore nel messaggio, per
          esempio il pid e il nome dell'host in testa a ogni riga.
\end{itemize}
Con due macchine, l'indirizzo fra parentesi quadre distingue i peer remoti; se non arriva niente, il colpevole tipico è
l'interfaccia scelta da \texttt{INADDR\_ANY} per l'iscrizione, o una VM in NAT.
\end{esercizio}

\begin{esercizio}{Iterazione 4: una chat guidata dagli eventi}
Si riparte dalla chat con i thread: funziona, e un client che smette di leggere congela tutti, perché il
\texttt{broadcast} si blocca con il lock della lista in mano.
\begin{itemize}
    \item \textbf{Che cosa cambia}: un thread solo. Niente \texttt{pthread\_create}, niente mutex: la lista è un array
          qualsiasi, e \texttt{accept} avviene quando il socket di ascolto è leggibile.
    \item \textbf{Che cosa è nuovo}: niente può bloccarsi. I descriptor diventano non bloccanti, e ogni client riceve una
          coda in uscita che entra nell'insieme di scrittura solo quando contiene qualcosa.
\end{itemize}
Misurare la memoria con 200 client connessi, contro la versione con i thread. Poi sospendere un client con Ctrl-Z: che
cosa succede agli altri 199? E che cosa fa la memoria del server mentre la coda di quel client cresce?

\svolgimento
Ogni client diventa una struttura con le due metà dello stato che prima stavano sullo stack del suo thread:

\begin{lstlisting}[language={[POSIX]C}, style=wnumbers, title={chatd.c: lo stato di un client}]
struct client {
    int    fd;                               /* -1 = posto libero */
    char   in[512];   size_t inlen;          /* la riga in costruzione */
    char   out[4096]; size_t outlen;         /* byte che aspettano di partire */
};
static struct client cl[MAXC];
static fd_set allset, wall;
\end{lstlisting}

\begin{lstlisting}[language={[POSIX]C}, style=wnumbers, title={chatd.c: il ciclo}]
for (;;) {
    fd_set rset = allset, wset = wall;
    if (select(maxfd + 1, &rset, &wset, NULL, NULL) < 0) { if (errno == EINTR) continue; break; }

    if (FD_ISSET(listen_fd, &rset)) accept_one();    /* nuovo client, O_NONBLOCK */
    for (int i = 0; i < MAXC; i++) {
        if (cl[i].fd != -1 && FD_ISSET(cl[i].fd, &rset)) on_readable(i);
        if (cl[i].fd != -1 && FD_ISSET(cl[i].fd, &wset)) on_writable(i);
    }
}
\end{lstlisting}

\texttt{on\_readable} legge quello che c'è, lo aggiunge a \texttt{in} e, per ogni \texttt{\textbackslash n} trovato,
chiama \texttt{queue(j, ...)} per ogni altro client $j$; su 0 o su un errore diverso da \texttt{EAGAIN} chiama
\texttt{drop(i)}. \texttt{on\_writable} è il pezzo «client lento» del capitolo. \texttt{drop} toglie il descriptor da
entrambi gli insiemi, lo chiude e libera il posto. Il secondo \texttt{if} ricontrolla \texttt{fd}, perché il primo può
aver appena scartato quel client.

Con un client sospeso, gli altri 199 \textbf{continuano a chattare}. Le righe per lui riempiono prima i buffer del kernel
(la sua finestra di ricezione e il nostro buffer di invio, centinaia di KB in tutto), poi la sua coda di 4096 byte; a quel
punto \texttt{queue} lo scarta. La memoria del server \textbf{non cresce}: le code sono di dimensione fissa. Con una coda
che si allarga con \texttt{realloc} crescerebbe senza limite, un messaggio alla volta, ed è proprio ciò che il limite
evita.

Quanto alla memoria: con il ciclo gli stati di 200 client sono 200 strutture da circa 4,6 KB in un solo processo; la
versione con i thread ha in più 200 stack, 200 thread nel kernel, e i cambi di contesto fra di loro.
\end{esercizio}

\begin{esercizio}{Iterazione 5: un pool dietro il ciclo}
Si riparte dall'iterazione 4: un solo flusso di controllo per tutti. Un client che fa lavorare il server (una ricerca, un
file letto dal disco) ferma ogni altra conversazione per tutto il tempo che serve.
\begin{itemize}
    \item \textbf{Che cosa cambia}: il ciclo smette di servire. Quando un descriptor è leggibile, esce dall'insieme
          osservato e va nella coda; un worker lo legge, fa il lavoro, accoda la risposta e restituisce il descriptor.
    \item \textbf{Perché non un pool di connessioni}: perché una connessione di chat non finisce mai, e sedici worker
          servirebbero sedici client per sempre. Il lavoro è il messaggio.
\end{itemize}
Misurare, con otto worker e un client che chiede qualcosa di lento, il ritardo visto dagli altri. Perché il descriptor deve
uscire dall'insieme osservato mentre lo tiene un worker? E dopo quanti worker la risposta smette di migliorare?

\svolgimento
Il ciclo e i worker si passano il descriptor in due direzioni: verso i worker con la coda limitata, all'indietro con una
\textbf{pipe} che il ciclo osserva insieme ai socket. È l'unico modo di svegliare un ciclo fermo in \texttt{select}.

\begin{lstlisting}[language={[POSIX]C}, style=wnumbers, title={chatd.c: andata e ritorno di un descriptor}]
/* nel ciclo: un client leggibile va a un worker */
if (FD_ISSET(fd, &rset)) {
    FD_CLR(fd, &allset);                     /* finche' lo tiene un worker, select non lo guarda */
    put_fd(&jobs, fd);
}

/* il worker */
for (;;) {
    int fd = get_fd(&jobs);
    handle(fd);                              /* legge, lavora, accoda la risposta */
    write(done[1], &fd, sizeof fd);          /* lo restituisce: 4 byte, scrittura atomica */
}

/* nel ciclo: done[0] e' nell'insieme osservato */
if (FD_ISSET(done[0], &rset)) {
    int fd;
    read(done[0], &fd, sizeof fd);
    FD_SET(fd, &allset);                     /* di nuovo osservato */
}
\end{lstlisting}

Ora le code dei client sono toccate dal ciclo e dai worker insieme, quindi tornano un mutex (uno per client basta).

Il descriptor esce dall'insieme perché, finché il worker non ha letto, il socket resta \textbf{leggibile}: \texttt{select}
tornerebbe subito a ogni giro, facendo girare il ciclo a vuoto, e il ciclo potrebbe consegnare lo stesso descriptor a un
secondo worker. Due worker sulla stessa connessione si dividerebbero le righe, e le risposte arriverebbero mescolate.

Con un solo client lento, il ritardo degli altri crolla già con \textbf{due} worker: uno resta occupato con la richiesta
lenta, l'altro smaltisce tutte le veloci. In generale si smette di migliorare quando i worker sono tanti quante le
richieste lente contemporanee più uno; e se il lavoro calcola invece di aspettare il disco, quando sono tanti quanti i
core.
\end{esercizio}

\Needspace{20\baselineskip}
% ══════════════════════════════════════════════════════════════
\section{Riepilogo}
% ══════════════════════════════════════════════════════════════

\riepilogo{
\begin{itemize}
    \item \texttt{select} si blocca su un insieme invece che su un descriptor, e restituisce quelli che si possono servire
          senza aspettare: un flusso di controllo per tutte le connessioni.
    \item Gli insiemi sono argomenti e risultati, quindi l'insieme principale si copia a ogni giro; e niente, sulla via da
          \texttt{select} a \texttt{select}, può bloccarsi.
    \item È questo che finalmente risponde al client lento: una \textbf{coda in uscita} per connessione, svuotata quando il
          socket dice di essere scrivibile.
    \item Un \textbf{pool} paga i suoi thread una volta e li tiene: il ciclo di \texttt{accept} produce, i worker
          consumano, e una coda limitata trasforma il sovraccarico in rifiuto invece che in esaurimento.
    \item Un pool serve tante connessioni lunghe quanti worker ha: quindi il lavoro è un messaggio, e qualcun altro si
          accorge che un messaggio è arrivato.
    \item Il \textbf{broadcast} raggiunge ogni host della sottorete e si ferma lì; il \textbf{multicast} raggiunge chi si è
          iscritto a un gruppo, su UDP, con una copia sul filo e nessuna idea di chi ascolta.
\end{itemize}
}
